% !TEX root = chapter-standalone.tex

\section{Upper and lower bounds}
\linkvideo{spring2021-tradeoffs:tradeoffs:orders:up-low-bounds} % Upper and lower bounds

\todotext{generalize upper and lower bound definitions to preorders}

\begin{ctdefinition}[Upper bounds in a poset]
    \label{def:least-upper-bound}
    \label{def:upper-bounds}
    Let $\posA$ be a poset and let $\subA \setsubseteq \posAset$ be a subset.
    An \emph{upper bound} of $\subA$ is:

    \constit
    \begin{enumerate}
        \item An element $\elb \setin \posAset$.
    \end{enumerate}

    \condit
    \begin{enumerate}
        \item $\elb$ dominates all elements of $\subA$: $\ela \posAleq \elb$ for all $\ela \setin \subA$.
    \end{enumerate}
    The set of all \maindef{upper bounds} of $\subA$ is denoted
    \begin{equation}
        \upperb \subA\definedas \makeset{ \elb \setin \posAset \mid \forall \ela \setin \subA \colon \ela \posAleq \elb }.
    \end{equation}
\end{ctdefinition}

\begin{ctdefinition}[Least upper bound / \SY{join} / supremum]\label{def:supremum}
    Let $\posA$ be a poset and let $\subA \setsubseteq \posAset$ be a subset.
    A \maindef{least upper bound} of $\subA$ is:

    \constit
    \begin{enumerate}
        \item An element $\elc \setin \posAset$.
    \end{enumerate}

    \condit
    \begin{enumerate}
        \item $\elc$ is an upper bound of $\subA$: $\elc \setin \upperb \subA$;
        \item $\elc$ is below every upper bound of $\subA$: $\elc \posAleq \elb$ for all $\elb \setin \upperb \subA$.
    \end{enumerate}
    If it exists, it is denoted $\join \subA$ or $\Sup\subA$, and also called the \maindef{join} or \maindef{supremum} of $\subA$.
\end{ctdefinition}

So, given~$\subA \setsubseteq \posAset$ and~$\elb \setin \posAset$,~$\elb = \join \subA$ if and only if
\begin{enumerate}
    \item $\ela \posAleq \elb, \ \forall \ela \setin \subA$, and
    \item $\ela \posAleq \elc, \ \forall \ela \setin \subA \Imp \elb \posAleq \elc.
          $
\end{enumerate}

If a least upper bound of a subset~$\subA \setsubseteq \posAset$ exists, it is unique (can you prove this?), so we speak of ``the'' least upper bound.

\begin{exercise}
    Let~\posA be a \SY{poset} and~$\subA \setsubseteq \posAset$ a subset of the underlying set of~\posA.
    Show that if~$\join \subA$ exists, then it is unique.
    For this, assume that~$\elb$ and~$\elc$ are both least upper bounds of~$\subA$, and then show that this assumption implies that in fact~$\elb = \elc$.
\end{exercise}
\begin{solution}
    Assume that~$\elb$ and~$\elc$ are both least upper bounds of~$\subA \setsubseteq \posA$.
    In other words, one knows~$\ela \posAleq \elb$ and~$\ela\posAleq \elc$ for all~$\ela\setin \subA$.
    However, one also has~$\elb\posAleq \elc$ and~$\elc\posAleq \elb$ (from~$\elb,\elc$ assumed to be both least upper bounds).
    Because of antisymmetry, this implies~$\elb=\elc$ and proves the uniqueness of least upper bounds in a poset.
\end{solution}

\begin{example}
    Consider the poset~\posA and its subset~$\subA$ depicted in \cref{fig:upper_bound_example_bis}.
    The \textcolor{red}{red} markers~$\textcolor{red}{\marker}$ represent the upper bound of~$\subA$.
    For this specific case, there is \emph{a single} least upper bound.
\end{example}

\begin{figure}[h!]
    \middlesag{upper_bound_2}
    \caption{Example of upper bounds and least upper bound for~$\subA$.}
    \label{fig:upper_bound_example_bis}
\end{figure}
\begin{example}
    Least upper bounds need not necessarily exist even in total orders.
    For instance, the subset
    \begin{equation}
        \posReals = \makeset{\ela \setin \reals \colon \ela>0}
    \end{equation}
    of the poset~\reals (with the usual ordering) does not have a least upper bound.
\end{example}

Analogously to the case of (least) upper bounds, we can define lower bounds and greatest lower bounds.

\begin{ctdefinition}[Lower bounds in a poset]
    \label{def:greatest-lower-bound}
    \label{def:lower-bounds}
    Let $\posA$ be a poset and let $\subA \setsubseteq \posAset$ be a subset.
    A \emph{lower bound} of $\subA$ is:

    \constit
    \begin{enumerate}
        \item An element $\elb \setin \posAset$.
    \end{enumerate}

    \condit
    \begin{enumerate}
        \item $\elb$ is dominated by all elements of $\subA$: $\elb \posAleq \ela$ for all $\ela \setin \subA$.
    \end{enumerate}
    The set of all \maindef{lower bounds} of $\subA$ is denoted
    \begin{equation}
        \lowerb \subA\definedas \makeset{ \elb \setin \posAset \mid \forall \ela \setin \subA \colon \elb \posAleq \ela }.
    \end{equation}
\end{ctdefinition}

\begin{ctdefinition}[Greatest lower bound / \SY{meet} / infimum]\label{def:infimum}
    Let $\posA$ be a poset and let $\subA \setsubseteq \posAset$ be a subset.
    A \maindef{greatest lower bound} of $\subA$ is:

    \constit
    \begin{enumerate}
        \item An element $\elc \setin \posAset$.
    \end{enumerate}

    \condit
    \begin{enumerate}
        \item $\elc$ is a lower bound of $\subA$: $\elc \setin \lowerb \subA$;
        \item $\elc$ is above every lower bound of $\subA$: $\elb \posAleq \elc$ for all $\elb \setin \lowerb \subA$.
    \end{enumerate}
    If it exists, it is denoted $\meet \subA$ or $\Inf \subA$, and also called the \maindef{meet} or \maindef{infimum} of $\subA$.
\end{ctdefinition}

\begin{exercise}
    Come up with an example of a subset~$\subA$ of a poset~\posA which has lower bounds but no greatest lower bound.
    Then, modify it to have a greatest lower bound.
\end{exercise}

\begin{solution}
    \begin{marginfigure}
        \centering
        \includesag{lower_bound_1}
        \caption{
            Example of lower bounds of~$\subA$.
        }
        \label{fig:lower_bound_1}
    \end{marginfigure}

    \begin{marginfigure}
        \centering
        \includesag{lower_bound_2}
        \caption{
            Example of lower bounds and greatest lower bounds of~$\subA$.
        }
        \label{fig:lower_bound_2}
    \end{marginfigure}

    In \cref{fig:lower_bound_1} you find an example of a subset~$\subA$ of a poset~\posA which has incomparable lower bounds.
    In \cref{fig:lower_bound_2} instead, there is a greatest lower bound.
\end{solution}
\linkvideo{spring2021-tradeoffs:tradeoffs:orders:top-bottom} % Top and bottom

\begin{ctdefinition}[Top and bottom]
    \label{def:top}
    \label{def:bot}
    Let $\posA$ be a \SY{poset}. A \emph{top} element of $\posA$ is:

    \constit
    \begin{enumerate}
        \item An element $\postop \setin \posAset$.
    \end{enumerate}

    \condit
    \begin{enumerate}
        \item $\postop$ is a least upper bound of the entire \SY{poset} $\posA$; equivalently, $\ela \posAleq \postop$ for all $\ela \setin \posAset$.
    \end{enumerate}
    Dually, a \emph{bottom} element $\posbot$ of $\posA$ is a greatest lower bound of the entire \SY{poset}; equivalently, $\posbot \posAleq \ela$ for all $\ela \setin \posAset$.
\end{ctdefinition}

\todotext{Define the notion of lattice and complete lattice here}

\begin{ctdefinition}[Power set as lattice]
    \label{def:power-set-as-lattice}
    Given a set $\setC$, its power set $\powerset \setC$ (the set of all subsets) is a \SY{lattice} given by:

    \constit
    \begin{enumerate}
        \item The carrier set $\powerset \setC$;
        \item The order given by inclusion:
              \begin{equation}
                  \setA\posleq \setB \definedas \setA\setsubseteq \setB;
              \end{equation}
        \item The \SY{join} given by the \SY{union of sets}:
              \begin{equation}
                  \setA\join \setB \definedas \setA\setunion \setB;
              \end{equation}
        \item The \SY{meet} given by the \SY{intersection of sets}:
              \begin{equation}
                  \setA\meet \setB \definedas \setA\setintersection \setB;
              \end{equation}
        \item The \SY{top element}, which is the set $\setC$ itself:
              \begin{equation}
                  \postop = \setC;
              \end{equation}
        \item The \SY{bottom element}, which is the empty set:
              \begin{equation}
                  \posbot = \Emptyset.
              \end{equation}
    \end{enumerate}
\end{ctdefinition}

\begin{marginfigure}
    \centering
    \includesag{060_powerset_coprod}
    \caption{}
    \label{fig:prod_coprod_power}
\end{marginfigure}

The Hasse diagram reported in \cref{fig:prod_coprod_power} illustrates the structure of the power set \SY{lattice} for three sets~$\setA,\setB,\setC\setin \powerset \stylesets{S}$.


\subsection{Minimal and maximal elements}

\todotext{Min and max are a thing that is here for DP; we don't generalize it to preorders and we need to think about where to put this material}

You know already the operators~$\min$/$\max$ that give the minimum/maximum values of a set of real numbers.
If the set is finite, the minimum and maximum always exist.
But for infinite sets, the minimum and maximum might not exist.
For example, consider the set of real numbers contained between 0 and 1, excluding the boundaries:
%
\begin{equation}
    \setA = \makeset{ \ela \setin \reals \colon 0 < \ela < 1 }.
\end{equation}
%
This set does not have a minimum or a maximum.

For a total order, if the minimum and maximum exist, then they are unique.
In a partial order, this is not the case.
We introduce the operators~$\Min$ and~$\Max$ that are the generalization to partial orders of~$\min/\max$.

\begin{ctdefinition}[Minimal elements]
    \label{def:Min}
    Let $\posA$ be a \SY{poset}. The \emph{minimal elements} function is:

    \constit
    \begin{enumerate}
        \item The function $\Min \colon \posPA \to \posAA$ that sends a subset $\subA$ to the minimal elements of that subset:
              \begin{equation}
                  \defmapperiodset{\Min}{\posPA}{\posAA}{\subA}{\makesett{
                          \setCel\setin \subA\colon
                          \vmiddle{
                              \prftree{
                                  \setDel\setin \subA
                              }{
                                  \setDel\posAleq \setCel
                              }{
                                  \setCel=\setDel
                              }
                          }
                      }}
              \end{equation}
    \end{enumerate}
    Note that $\Min(\subA)$ could be empty.
\end{ctdefinition}

\begin{ctdefinition}[Maximal elements]
    \label{def:Max}
    Let $\posA$ be a \SY{poset}. The \emph{maximal elements} function is:

    \constit
    \begin{enumerate}
        \item The function $\Max \colon \posPA \to \posAA$ that sends a subset $\subA$ to the maximal elements of that subset:
              \begin{equation}
                  \defmapperiodset{
                      \Max
                  }{
                      \posPA
                  }{
                      \posAA
                  }{
                      \subA
                  }{
                      \makesett{
                          \setCel\setin \subA\colon
                          \vmiddle{
                              \prftree{
                                  \setDel\setin \subA
                              }{
                                  \setDel\posAgeq \setCel
                              }{
                                  \setCel=\setDel
                              }
                          }
                      }
                  }
              \end{equation}
    \end{enumerate}
    Note that $\Max(\subA)$ could be empty.
\end{ctdefinition}

